\section{Related Work}
\label{sec:related}

\para{Activation steering.}
\citet{turner2023activation} introduce activation addition (\actadd), computing steering vectors as mean activation differences between contrastive prompt pairs and adding them to the residual stream. \citet{zou2023representation} extend this idea under the umbrella of representation engineering, using linear probes and interventions for top-down model control. Several works have refined these methods: \citet{rimsky2024steering} scale steering to safety-relevant behaviors, and more recent work explores conditional activation steering \citep{lee2024programming}, multi-property steering \citep{wang2024multi}, and angular steering \citep{wang2025angular}. A key finding from \citet{tan2024analyzing} is that steering vectors have substantial reliability limitations---their effectiveness varies across inputs and can reflect spurious correlations. All of these methods operate in the residual stream; we extend the comparison to the MLP intermediate space.

\para{MLP interpretability.}
The MLP sublayer performs a large fraction of a transformer's computation, yet its intermediate space remains less studied than the residual stream. \citet{cunningham2023sparse} train sparse autoencoders (SAEs) on model activations and find that the learned features are more interpretable than individual neurons. \citet{dunefsky2024transcoders} introduce transcoders---wide, sparse approximations of MLP sublayers that map input to output directly---enabling input-invariant circuit analysis through the MLP. Most relevant to our work, \citet{shafran2025decomposing} decompose MLP intermediate activations using semi-nonnegative matrix factorization (SNMF) and show that the resulting features outperform both SAEs and supervised baselines (DiffMeans) at causal steering. Their features reveal hierarchical, compositional structure where related concepts share core neurons. \citet{marks2026language} provide further evidence that MLP computations are sparse in the neuron basis, with only a small fraction contributing to any given task.

\para{Linear representations in neural networks.}
The linear representation hypothesis \citep{nanda2023emergent, park2024linear} posits that high-level concepts are encoded as linear directions in activation space. This hypothesis underpins both linear probing and activation steering. If concepts are indeed linear, then mean-difference vectors should align with concept directions, and wider spaces with less superposition should be easier to steer. We test this prediction by comparing linear probe accuracy across spaces.

\para{Our contribution.}
Prior work on MLP-space steering uses learned features (transcoders, SNMF) rather than direct mean-difference vectors. We ask the simpler question: does the basic mean-difference approach---the workhorse of activation steering---work better or worse in the MLP intermediate space? This comparison isolates the effect of the intervention space from the effect of the feature extraction method.
